% From power-on to the game: the numbered list of screens, in order.
\begin{enumerate}
\item \textbf{Message on the MSX screen.} For about two seconds the MSX's own screen shows \texttt{CAPCOM 1985},
\texttt{DIHALT STUDIO 2026}, \texttt{VERSION \portVersion} and \texttt{V9990 found!}. If it says \texttt{THIS GAME RUNS ONLY ON MSX TURBO R} (the computer is
not an MSX turbo R) or \texttt{V9990 NOT found!!!} (no V9990), the game stops there.
\item \textbf{DiHalt logo.} With a MoonSound or Wozblaster, the logo stays on the V9990 screen throughout the ten-second DiHalt Studio fanfare. Other configurations show the logo for four seconds in silence.
\item \textbf{Cover.} The full-screen illustration. It waits until you press \key{SPACE} or the joystick fire button.
\item \textbf{Homage message.} A short text is typed on the screen: this game is a homage to the original Commando by
Capcom.
\item \textbf{Title.} The screen with the bonus scores and the blinking \texttt{PRESS FIRE TO START}, then the big
\textsc{Commando} logo with the ranking of the best seven scores. It repeats until you start.
\item \textbf{Start.} Press \key{1} or fire for a normal game, or \key{2} or the grenade button for an easy game (the
enemy bullets are slower). There are no coins and no credits, and the game is for one player.
\item \textbf{The game.} A helicopter drops you at the start of area~1 and the game begins.
\end{enumerate}

When your last life is lost the screen shows \texttt{GAME OVER}. If your score is among the best seven you can write your
name (section~8); then the ranking is shown and the title screens start again.
